\chapter{Tanggung Jawab Kita}\label{ch:g9:our-responsibility}

Sembilan tahun buku ini berakhir di sini, dan benangnya
bertemu pada sebuah kenyataan yang sederhana: \omterm{def:g5:biodiversity:species}{spesies} yang
dapat membaca halaman ini adalah satu-satunya \omterm{def:g5:biodiversity:species}{spesies} yang
memahami permesinan tempatnya berada --- dan pemahaman,
seperti pertama kali dikatakan \cref{ch:g8:contraception},
membawa serta pilihan. Bab penutup ini mengumpulkan tanggung
jawab yang diciptakan pemahaman itu: terhadap kesehatan diri
sendiri, terhadap kesehatan orang lain, dan terhadap jaring
\omterm{def:g6:exploring-our-environment:components}{makhluk hidup} yang mengangkut kita semua.

\section{Tanggung jawab atas permesinan sendiri}

\begin{proposition}[Kesehatan sebagai pengelolaan diri yang berpengetahuan]\label{prop:g9:our-responsibility:self}
Benang kesehatan buku ini, disimpulkan:
\begin{itemize}
  \item kontrak hariannya
        (\cref{prop:g5:healthy-choices:contract}) ---
        makanan, gerak, \omterm{prop:g1:healthy-habits:needs}{tidur}, kebersihan --- adalah
        perawatan organ yang kini dapat kamu beri alasan
        sistem demi sistem
        (\cref{prop:g7:heart-and-circulation:keep},
        \cref{met:g7:lungs-and-gas-exchange:keep},
        \cref{met:g8:protecting-the-nervous-system:manual});
  \item mekanisme zatnya sudah diketahui
        (\cref{prop:g8:protecting-the-nervous-system:substances}),
        ketaksetangkupan jebakannya sudah dihargai
        (\cref{rem:g5:healthy-choices:never});
  \item pilihan reproduksinya punya petanya dan alamatnya
        (\cref{ch:g8:contraception});
  \item dan riwayat tubuhnya sebagian dibagikan
        (\cref{ch:g9:heredity-and-traits}): mengetahui
        kecenderungan sebuah keluarga --- ubannya yang dini,
        \omterm{def:g4:heart-and-blood:heart}{jantungnya} yang berkerak --- bukanlah takdir
        melainkan sebuah jadwal perawatan
        (\cref{met:g9:unique-individuals:claims}: rentang,
        yang dijalani).
\end{itemize}
\end{proposition}
\admitted

\begin{example}[Pemeriksaan berkala, dibaca sebagai biologi]\label{ex:g9:our-responsibility:checkup}
Sebuah pemeriksaan berkala yang biasa adalah buku ini di
dalam satu janji temu: \omterm{prop:g4:heart-and-blood:pulse}{denyut nadi} dan tekanan \omterm{prop:g4:journey-of-food:absorb}{darahnya}
(jaringannya dibaca dari tepi jalan,
\cref{ex:g7:heart-and-circulation:measures}), palu
refleksnya (\cref{pb:g8:nervous-communication:1}), buku
besar \omterm{prop:g9:immune-defenses:vaccination}{vaksinasinya} (\cref{pb:g9:immune-defenses:1}), riwayat
keluarganya (kecenderungan yang dibagikan), dan nasihat
tetapnya --- kontraknya, sekali lagi. Kedokteran sebagian
besar adalah perawatan terapan atas mekanisme yang buku
manualnya kini kamu miliki.
\end{example}

\section{Tanggung jawab terhadap orang lain}

\begin{proposition}[Kesehatan adalah permesinan bersama]\label{prop:g9:our-responsibility:others}
Tiga mekanisme menjadikan kesehatan sebuah rekening bersama:
\begin{itemize}
  \item \emph{penularannya}
        (\cref{def:g9:microbes-and-infection:stages}):
        \omterm{def:g1:healthy-habits:hygiene}{kebersihan diri} yang kamu jaga, tertutupnya batukmu dan
        hari sakitmu di rumah adalah perlindungan bagi orang lain;
  \item \emph{\omterm{def:g9:immune-defenses:memory}{kekebalan} yang berkumpul}
        (\cref{prop:g9:immune-defenses:vaccination}):
        \omterm{prop:g9:immune-defenses:vaccination}{vaksinasi} menjaga, lewat dirimu, mereka yang tak
        dapat dilatih;
  \item \emph{harta bersama kedokterannya}: antibiotik yang
        dijaga tetap bekerja
        (\cref{exo:g9:how-species-change:10}) dan \omterm{prop:g4:journey-of-food:absorb}{darah} yang
        disumbangkan (\cref{ex:g9:unique-individuals:medicine}
        --- daftarnya berjalan dengan pendonor) adalah sumber
        daya bersama yang dibelanjakan atau diisi oleh
        pilihan perorangan.
\end{itemize}
\end{proposition}
\admitted

\begin{example}[Rekening bersama seorang remaja]\label{ex:g9:our-responsibility:account}
Catatan biasa sepekan itu: tinggal di rumah karena demam
(penularannya diputus); suntikan ulangnya mutakhir (kolamnya
tertahan); sebuah kursus antibiotik musim semi yang
dituntaskan (harta bersamanya tak terbelanjakan); sebuah
donor \omterm{prop:g4:journey-of-food:absorb}{darah} pertama yang direncanakan pada usia delapan
belas (daftarnya terisi). Tak satu pun kepahlawanan ---
pelajaran \cref{ex:g5:healthy-choices:day} pada skala sebuah
kota: kesehatan masyarakat adalah hari-hari biasa, yang dijaga.
\end{example}

\section{Tanggung jawab terhadap dunia yang hidup}

\begin{proposition}[Penatalayanannya, dinyatakan ulang beserta semua alasannya]\label{prop:g9:our-responsibility:stewardship}
Benang alamnya, disimpulkan beserta mekanismenya:
\begin{itemize}
  \item setiap \omterm{def:g3:food-chains:chain}{rantai makanan} bermula hijau
        (\cref{prop:g6:plants-are-producers:firstlink}),
        setiap meja makan adalah keluaran sebuah ekosistem
        (\cref{prop:g6:producing-our-food:costs});
  \item jaringnya menyerap guncangan lewat keragamannya
        (\cref{prop:g5:biodiversity:web}), dan keragamannya
        adalah cadangan yang dijalankan seleksinya
        (\cref{prop:g9:unique-individuals:reserve},
        \cref{exo:g9:how-species-change:12}) ---
        \omterm{def:g5:biodiversity:biodiversity}{keanekaragaman hayati} bukan pemandangan melainkan
        modal kerja;
  \item kepunahan menutup pintunya selamanya
        (\cref{prop:g5:biodiversity:loss}) --- satu-satunya
        pengecualian yang disengaja, yaitu cacar,
        membuktikan betapa mutlaknya kata itu
        (\cref{exo:g9:immune-defenses:12});
  \item aturan pemakaian yang lestari sudah diketahui
        (\cref{prop:g5:people-and-nature:lasting}),
        pelanggarannya dapat didiagnosis
        (\cref{ex:g7:oxygen-in-the-water:waste}),
        perbaikannya sudah diperagakan
        (\cref{ex:g5:people-and-nature:river},
        \cref{ex:g5:biodiversity:storks}).
\end{itemize}
Penatalayanan adalah pola pikir perawatan yang sama dengan
kesehatan, diterapkan pada tubuh yang lebih besar tempat
kita tinggal.
\end{proposition}
\admitted

\begin{example}[Untuk apa biologi seorang warga]\label{ex:g9:our-responsibility:citizen}
Suara pemilih, pembelian dan kebiasaan kehidupan dewasa
terus berdatangan berpakaian sebagai pertanyaan yang sudah
dibekali buku ini: Apakah perikanan ini mengambil lebih
cepat daripada pembaruannya? Apakah kebijakan ini menjaga
\omterm{def:g2:where-animals-live:habitat}{habitat} penyerbuknya atau mengaspalnya? Apakah pola makan
ajaib, obat atau ``\omterm{def:g9:genes-and-dna:gene}{gen} untuk X'' yang diklaim ini runtut
(\cref{met:g9:unique-individuals:claims},
\cref{met:g8:contraception:evaluate})? Perlukah antibiotik
ini ada di dalam palung pakan ini? Seorang warga yang dapat
menjalankan audit itu adalah tujuan sembilan tahun babnya.
\end{example}

\begin{method}[Cara sang lulusan]\label{met:g9:our-responsibility:method}
Meninggalkan Buku 1, bawalah caranya lebih daripada faktanya:
\begin{enumerate}
  \item amati sebelum menyimpulkan
        (\cref{met:g1:animals-around-us:observe}
        memulainya; setiap sensus dan buku catatan sesudahnya);
  \item ujilah dengan adil --- satu perbedaan, sebuah kontrol
        (\cref{met:g4:what-plants-need:fairtest}, dipakai
        pada tiap tahun sesudahnya);
  \item sebutkan mekanismenya --- tidak ada klaim yang
        teraudit sampai mata rantainya disebutkan
        (\cref{met:g8:contraception:evaluate});
  \item timbanglah saksi yang saling bebas
        (\cref{met:g9:evidence-of-evolution:weigh});
  \item dan hormatilah permesinan yang hidup, di dalam
        tubuhmu dan di luarnya: ia tua, halus, dapat
        diperbaiki dalam batasnya --- dan milikmu untuk dijaga.
\end{enumerate}
\end{method}

\begin{remark}[Apa yang menyusul]\label{rem:g9:our-responsibility:next}
Jilid sekolah menengahnya mengangkat benang ini pada
perbesaran berikutnya: permesinan molekul \omterm{def:g5:first-look-at-cells:cell}{selnya}, genetika
beserta matematikanya yang lengkap, sistem \omterm{def:g9:immune-defenses:memory}{kekebalan} dan
\omterm{def:g8:nervous-communication:neuron}{sarafnya} dengan sungguh-sungguh, ekologi sebagai sebuah ilmu
kuantitatif, dan \omterm{prop:g9:how-species-change:selection}{evolusi} sebagai kerangka yang memang ia
adanya. Pertanyaannya akan sama, ditanyakan dari lebih
dekat. Bawalah caranya.
\end{remark}

\section{Latihan}

\begin{exercise}[$\star$]\label{exo:g9:our-responsibility:1}
Sebutkan keempat untai tanggung jawab diri pada
\cref{prop:g9:our-responsibility:self}.
\end{exercise}

\begin{exercise}[$\star$]\label{exo:g9:our-responsibility:2}
Sebutkan ketiga mekanisme yang menjadikan kesehatan
bersama, masing-masing dengan babnya.
\end{exercise}

\begin{exercise}[$\star$]\label{exo:g9:our-responsibility:3}
Mengapa kecenderungan keluarga yang diketahui adalah sebuah
jadwal alih-alih sebuah takdir?
\end{exercise}

\begin{exercise}[$\star$]\label{exo:g9:our-responsibility:4}
Bacalah tiga alat pemeriksaan berkalanya sebagai babnya.
\end{exercise}

\begin{exercise}[$\star$]\label{exo:g9:our-responsibility:5}
Mengapa \omterm{def:g5:biodiversity:biodiversity}{keanekaragaman hayati} adalah ``modal kerja''
alih-alih pemandangan? Dua mekanisme.
\end{exercise}

\begin{exercise}[$\star$]\label{exo:g9:our-responsibility:6}
Nyatakan kelima langkah cara sang lulusan dari ingatan.
\end{exercise}

\begin{exercise}[$\star\star$]\label{exo:g9:our-responsibility:7}
Berikan alasan bagi hari sakit di rumah sebagai perbuatan
kesehatan masyarakat, dengan tahapnya disebutkan.
\end{exercise}

\begin{exercise}[$\star\star$]\label{exo:g9:our-responsibility:8}
Vaksinasimu adalah perisai siapa? Jawablah dengan kolamnya,
lalu sebutkan dua yang bergantung padanya.
\end{exercise}

\begin{exercise}[$\star\star$]\label{exo:g9:our-responsibility:9}
Auditlah sebagaimana \cref{ex:g9:our-responsibility:citizen}
akan mengauditnya: ``semprotan baru kami membunuh setiap
\omterm{prop:g3:sorting-living-things:groups}{serangga} di kebun \omterm{prop:g3:flowers-fruits-seeds:transform}{buahnya} --- perlindungan total.''
\end{exercise}

\begin{exercise}[$\star\star$]\label{exo:g9:our-responsibility:10}
Jelaskan ``antibiotik adalah sebuah harta bersama'' dengan
seleksi bangsalnya --- dan dua kebiasaan yang menjaga harta
bersamanya tetap terisi.
\end{exercise}

\begin{exercise}[$\star\star$]\label{exo:g9:our-responsibility:11}
Terapkan kelima langkah caranya pada sebuah unggahan ``obat
ajaib'' yang viral, satu langkah pada satu waktu.
\end{exercise}

\begin{exercise}[$\star\star\star$]\label{exo:g9:our-responsibility:12}
Tuliskan sendiri alinea penutup bukunya: satu kalimat
masing-masing untuk \omterm{def:g5:first-look-at-cells:cell}{selnya}, tubuhnya, jaringnya, pohon
keluarganya --- dan tempat sang pembaca pada keempatnya.
\end{exercise}

\section{Soal: Audit Kelulusannya}

\begin{problem}[{Soal akhir pekan --- setahun sebuah kota,
diaudit dengan seluruh bukunya}]\label{pb:g9:our-responsibility:1}
Setahun kotamu (rekaan) dalam tajuk berita: sebuah gelombang
flu musim dingin; sebuah perdebatan musim semi tentang
menyemprot pohon planenya; sebuah kematian massal \omterm{prop:g3:sorting-living-things:groups}{ikan} musim
panas di kolam tamannya; sebuah kampanye musim gugur untuk
cagar rawa yang baru; dan, sepanjang tahun, gerakan
\omterm{prop:g9:immune-defenses:vaccination}{vaksinasi} dan donor kliniknya.

\medskip\noindent\textbf{Bagian I --- Berkas kesehatannya.}
\begin{enumerate}
  \item Gelombang flunya: sebutkan jalur penularan yang
        dibidik nasihat kotanya, dan mengapa kliniknya
        memvaksinasi staf panti jompo lebih dahulu.
  \item Gerakannya: nyatakan apa yang diisi daftar \omterm{prop:g4:journey-of-food:absorb}{darah} dan
        buku besar \omterm{prop:g9:immune-defenses:vaccination}{vaksinasinya} masing-masing, dan siapa
        yang menariknya.
  \item Seorang kolumnis mengolok-olok ``keributan soal cuci
        tangan''. Jawablah dengan bangsal pada
        \cref{ex:g9:microbes-and-infection:history} --- satu
        alinea.
  \item Apoteknya melaporkan kursus antibiotik yang tak
        tuntas berdatangan kembali. Harta bersama yang mana
        yang dibelanjakan, oleh mekanisme yang mana?
\end{enumerate}

\medskip\noindent\textbf{Bagian II --- Berkas alamnya.}
\begin{enumerate}[resume]
  \item Perdebatan penyemprotannya: jalankan audit kebun
        \omterm{prop:g3:flowers-fruits-seeds:transform}{buahnya} (\cref{exo:g9:our-responsibility:9}) pada
        ``bunuh setiap \omterm{prop:g3:sorting-living-things:groups}{serangga}'' --- siapa lagi yang mati,
        dan jasa cuma-cuma apa yang berhenti?
  \item Kematian massal \omterm{prop:g3:sorting-living-things:groups}{ikannya}: sepekan panas tanpa angin
        dan sebuah kolam hijau kacang. Diagnosislah dengan
        \cref{ch:g7:oxygen-in-the-water}, lalu berikan
        resepnya.
  \item Kampanye rawanya: belalah cagarnya dengan tiga
        mekanisme --- \omterm{def:g2:where-animals-live:habitat}{habitatnya}, jaringnya, cadangan
        keragamannya --- satu kalimat masing-masing.
  \item Seorang pemilik \omterm{def:g6:decomposers-and-soil:soil}{tanah} bertanya ``apa yang pernah
        dikerjakan rawanya bagi kotanya?'' Jawablah dengan
        buku besar kerja cuma-cumanya
        (\cref{exo:g5:biodiversity:11}, penyimpanan banjir
        sungainya pada \cref{ex:g5:people-and-nature:river}).
\end{enumerate}

\medskip\noindent\textbf{Bagian III --- Penutup sang
lulusan.}
\begin{enumerate}[resume]
  \item Pilihlah keputusan terbaik dan terburuk tahun itu
        menurut aturan pemakaian yang lestari, lalu berikan
        alasan bagi tiap vonisnya.
  \item Perlihatkan kelima langkah caranya sedang bekerja
        pada bagaimana kematian massal \omterm{prop:g3:sorting-living-things:groups}{ikannya} dipecahkan
        --- dari pengamatan pertamanya sampai mekanisme yang
        diterima.
  \item Susunlah tajuk rencana akhir tahun koran kotanya
        dalam empat kalimat: kesehatan yang dipikul bersama,
        alam yang disimpan, klaim yang diaudit, cara yang
        dijaga.
  \item Baris terakhir halaman terakhir Buku 1 adalah
        milikmu untuk ditulis: satu kalimat, yang ditujukan
        kepada pembaca yang memulai pada siput
        \cref{ch:g1:living-or-not}.
\end{enumerate}
\end{problem}
